% Generated draft: requires mathematical and source review.
\documentclass[11pt]{article}
\usepackage[margin=25mm]{geometry}
\usepackage{fontspec}
\setmainfont{Latin Modern Roman}
\usepackage{amsmath,amssymb,mathtools,unicode-math}
\setmathfont{Latin Modern Math}
\usepackage{enumitem,needspace,xurl,hyperref}
\setlength{\parindent}{0pt}
\setlength{\parskip}{5pt}
\setlength{\emergencystretch}{3em}
\newcommand{\N}{\mathbb{N}}
\newcommand{\Z}{\mathbb{Z}}
\newcommand{\Q}{\mathbb{Q}}
\newcommand{\R}{\mathbb{R}}
\newcommand{\C}{\mathbb{C}}
\newcommand{\F}{\mathbb{F}}
\newcommand{\A}{\mathbb{A}}
\newcommand{\PP}{\mathbb{P}}
\newcommand{\E}{\mathbb{E}}
\newcommand{\eps}{\varepsilon}
\newcommand{\dd}{\,\mathrm{d}}
\begin{document}
\section*{UnsolvedMath ID: 20000379}
\textbf{Source problem code:} catalogue entry
\subsection{Definitions and mathematical statement}
Brian Conrey's AIM question compares high-rank quadratic twists according to a local value. The supplied report explicitly reconstructs the setting as follows: fix odd $q$, an elliptic curve $E/\F_q(T)$, a finite rational good-reduction place $P=(T-a)$ with $a\in\F_q$, and monic squarefree polynomials $f$ of degree $n$. Let $E^f$ be the quadratic twist and let $h(f)$ denote the event $\operatorname{rank}E^f\geq2$. The original PDF does not state these field/place/family assumptions or define $\F_q[T]_n$; the rank convention also needs confirmation.

Partition the chosen family $\mathcal S_n$ into $\mathcal S_n^+$, $\mathcal S_n^-$ and $\mathcal S_n^0$ by $f(a)$ being a nonzero square, a nonsquare, or zero. The literally printed ratio is
\[
 R_n=\frac{\#\{f\in\mathcal S_n^+:h(f)\}}
 {\#\{f\in\mathcal S_n^-\cup\mathcal S_n^0:h(f)\}}.
\]
Indeed ``not a nonzero square'' includes zero, a ramified state. This differs from the cleaned unramified ratio with only $\mathcal S_n^-$ in the denominator. The extracted proposed local constant involves point counts on $E_P$ and its quadratic twist, but the radical/fraction formatting is damaged, so its exact value $C_P$ is not silently reconstructed. The source motivates the comparison by number-field numerical evidence, random matrix theory, moments of $L$-functions and BSD.

The report's unconditional local-counting theorem is
\[
 A_n:=\#\mathcal S_n^++\#\mathcal S_n^-=
 \frac{(q-1)(q^n-(-1)^n)}{q+1},
\]
\[
 \#\mathcal S_n^+-\#\mathcal S_n^-=
 \begin{cases}0&n\text{ odd},\\1-q&n>0\text{ even}.
 \end{cases}
\]
For fixed $q$ as $n\to\infty$, the resulting proportions in the full squarefree family are respectively
\[
 \frac{q}{2(q+1)},\qquad\frac{q}{2(q+1)},\qquad\frac1{q+1}.
\]
No high-rank asymptotic follows from these unconditional counts alone.

Put $r_\epsilon(n)=\#\{f\in\mathcal S_n^\epsilon:h(f)\}/\#\mathcal S_n^\epsilon$ for $\epsilon=+,-,0$. Conditionally, if the rates are defined, $r_-(n)>0$ eventually and
\[
 r_+(n)/r_-(n)\to C,\qquad r_0(n)/r_-(n)\to\lambda
\]
with finite nonnegative limits, the report gives
\[
 R_n\longrightarrow\frac{C}{1+2\lambda/q},
\]
whereas the cleaned square/nonsquare ratio tends to $C$. If $C=C_P$ and $\lambda=1$, the literal prediction is $qC_P/(q+2)$ rather than $C_P$. The existence and values of these rank-conditioned limits are unproved in the supplied record, which explicitly labels this consequence conditional.
\subsection{Short English statement}
Compare high-rank quadratic twists with square and nonsquare local values. The literal denominator also counts zeros, a positive-density ramified state. The report proves three-state local proportions, but its corrected limiting high-rank ratio remains conditional on unproved relative rank frequencies and on the explicit reconstructed family assumptions.
\subsection{Sources}
\begin{itemize}
\item Local registry: \url{data/unsolvedmath/problems.json}, permanent ID $20000379$, source code \texttt{catalogue entry}; damaged original question, supplied reconstruction and full research summary. Statement status reconstructed and unverified.
\item AIM workshop ``Arithmetic statistics over finite fields and function fields'', Question 1, attributed to Brian Conrey: \url{https://aimath.org/pastworkshops/arithstatffieldproblems.pdf}. Canonical location \texttt{aim-algebraic-number-theory-notes.json}, \texttt{notes[30]}; original remarks and literature field empty. No PDF was opened here.
\item Original motivation: number-field numerical evidence, random matrix theory, moments of $L$-functions and BSD; no full references supplied. The exact radical/point-count constant is not recoverable from the included extraction.
\item Supplied machine-generated research attempt 1: conditional result; problem open at the run and currently partially solved; rigor metadata ``conditional'', novelty confidence low. The local counting is reported unconditional, the rank-ratio consequence conditional. The non-peer-reviewed aid requires independent verification before citation.
\end{itemize}
\end{document}
